% Pista ana-23j2-q1 · Anàlisi
% Procedència: PAU Catalunya 2023 · Convocatòria de juny · Sèrie 5 · Qüestió 1
\documentclass[11pt,a4paper]{article}
\usepackage{pista}
\pistainfo{Catalunya 2023}{Convocatòria de juny}{Sèrie 5}

\begin{document}

\section*{\textcolor{blauPista}{Qüestió 1}}

\noindent Considereu les funcions $f(x) = -x^{2} + x + 6$ i $g(x) = -9x + 3x^{2}$.

\subsection*{\textit{a)} \normalfont Calculeu l'àrea de la regió delimitada per les dues funcions. \hfill\textnormal{[1,25~punts]}}

\pista{Primer necessites saber \emph{on} es tallen les dues corbes: iguala $f(x) = g(x)$ i resol l'equació de segon grau que en surt; obtindràs dues abscisses, que seran els límits d'integració. L'àrea entre dues corbes és $\int_{x_1}^{x_2} \bigl(\text{superior} - \text{inferior}\bigr)\,dx$, així que abans de res decideix quina funció queda per damunt en aquest interval (n'hi ha prou d'avaluar les dues en un punt intermedi). Compte amb els signes en restar $f - g$: et quedarà un sol polinomi de segon grau fàcil d'integrar amb la regla de Barrow.}

\subsection*{\textit{b)} \normalfont Trobeu l'equació de la recta tangent a la funció $f(x)$ en el punt $(-2, 0)$. Representeu aquesta recta tangent i les funcions $f(x)$ i $g(x)$ en uns mateixos eixos de coordenades. \hfill\textnormal{[1,25~punts]}}

\pista{El pendent de la tangent és $f'(-2)$. Com que el punt $(-2, 0)$ ja és a la recta, en tens prou per a determinar l'ordenada a l'origen. Per al dibuix, recorda que $f$ i $g$ són paràboles: calcula'n el vèrtex (on s'anul·la la derivada) i els talls amb l'eix $X$; la recta la traces amb dos punts, per exemple el de tangència i el tall amb l'eix $Y$.}

\end{document}
